% Appendix B — robustness and threats to validity (expanded from §Limitations)

This appendix expands the limitations listed in the main paper.

\paragraph{B.1 Selection into the repository universe.}
We restrict to \PLACEHOLDERPRACTIVE{} critical OS software repositories ranked by OpenSSF \texttt{criticality\_score} (\Cref{app:repo-selection}). The resulting sample has median stars \PLACEHOLDERSTARMED{} (range \PLACEHOLDERSTARMIN{}--\PLACEHOLDERSTARMAX{}) and a criticality-score range of \PLACEHOLDERSCOREMIN{}--\PLACEHOLDERSCOREMAX{} --- it skews toward load-bearing, multi-org, heavily-cross-referenced open-source infrastructure (language toolchains, runtimes, foundation projects) rather than uniform popularity. Effects we estimate apply to ``critical OS software repositories''. We do not generalise to GitHub as a whole. In particular, chain-length growth in lower-criticality repositories may be smaller or larger than what we report, and we do not have data to say which.

\paragraph{B.2 Detection false-negatives: ``undercover'' agents.}
Our event-level detection (\S\ref{sec:methods}, \Cref{app:allowlist}) relies on two signals: bot-account logins, and \texttt{Co-Authored-By} trailers naming an AI tool. Commercial coding agents have been reported to include configuration options that suppress attribution in outward artefacts~---~e.g., an ``undercover mode'' reportedly present in at least one major coding agent's system prompt, revealed via a source-code leak in early 2026. PRs produced by such configurations appear as ordinary human activity to our detector. Consequently the measured AI-agent participation share in Figure~\ref{fig:headline} is a \emph{lower bound}. This does not overturn the paper's main claim (participation is growing) but tightens its interpretation: the true level is above what we measure, by an unknown amount.

\paragraph{B.3 Thin AI-approval slice in the within-PR test.}
Among our \PLACEHOLDERNREPOS{} repositories and \PLACEHOLDERNPRS{} PRs, only \PLACEHOLDERWITHINN{} carry an approving review from an AI bot, and only \PLACEHOLDERWITHINAIN{} of those are also AI-authored. This is not a sampling accident: fewer than ten bots in our allowlist ever emit an ``APPROVED'' review state, and only two (Cubic, CodeRabbit) do so routinely. The two-proportion $z$-test on the 24 vs.\ 137 cells has limited power. We report it as a null on explicit self-preference, with the understanding that larger-sample follow-ups are needed to distinguish a small positive effect from zero at the $\pm 5$~pp level.

\paragraph{B.4 Merger attribution: authors and reviewers, not mergers.}
We initially considered ``AI merger'' as a third dimension (author $\times$ reviewer $\times$ merger). No merges in our universe are attributed to AI bot accounts on our allowlist, so the merger dimension collapses to ``human'' on every observed PR. This is a substantive limitation: a human clicks the merge button on every merged PR in our window, but AI can still act \emph{through} that human---most obviously when an AI review's ``approved'' state is what gates a maintainer's decision to merge. Our observational design cannot separate a maintainer who merged \emph{because of} an AI approval from one who merged in spite of it. We also caution that our detector excludes merges performed under a maintainer's personal-access token by an AI app or merge-queue automation acting on their behalf. Such merges would appear as human in both the REST \texttt{merged\_by} field and the GraphQL timeline actor, and therefore cannot be recovered from the public event log alone.

\paragraph{B.5 Event timestamp ordering and chain definition.}
AI$\to$AI chains are defined on events sorted by timestamp within a PR. GitHub timestamps are accurate to seconds. Ties are broken by the API's default order (commit $<$ review $<$ review-comment $<$ issue-comment). Alternative tie-breaking does not change the head of the distribution meaningfully (mean and $p_{95}$ shift by $<$0.1 events). The tail (max chain = 13) is insensitive.

\paragraph{B.6 Reproducibility.}
Every number in the paper is derived from one of \texttt{data/pr\_events.parquet} (one row per actor/event) or \texttt{data/chains.parquet} (one row per PR, chain aggregates). These are produced by \texttt{scripts/03\_classify\_prs.py} and \texttt{05\_compute\_chains.py} respectively. The within-PR two-proportion test is produced by \texttt{scripts/06c\_within\_pr.py}. The upstream PR-event data (\texttt{data/prs/*.jsonl}) is collected by \texttt{scripts/02\_fetch\_prs.py} against the \PLACEHOLDERNREPOS{} repositories in \texttt{data/repos.json}, which itself is produced by \texttt{scripts/01\_build\_repo\_list.py} (\Cref{app:repo-selection} details this step). Given a valid GitHub token, the full pipeline reproduces the paper's numbers up to GitHub's own indexing non-determinism.
